\subsection{Properties of product measures}

If \(\lambda\) (resp. \(\mu\) ) is a measure on X (resp. Y) and \(\nu\) is the product measure of \(\lambda\) by \(\mu\) , then the restriction of \(\nu\) to \(\mathcal{K}(X;C)\otimes_{\mathbf{C}}\mathcal{K}(Y;C)\) is none other than the tensor product \(\lambda\otimes\mu\) of the two linear forms \(\lambda\) and \(\mu\) (A, II, §3, No.~2), because the relation (1) of No.~1 may be written

\[
\langle g \otimes h, \nu \rangle = \langle g, \lambda \rangle \langle h, \mu \rangle = \langle g \otimes h, \lambda \otimes \mu \rangle
\]

for all \(g \in \mathcal{K}(X; \mathbf{C})\) and \(h \in \mathcal{K}(Y; \mathbf{C})\) . By an abuse of notation, we shall also denote by \(\lambda \otimes \mu\) the product measure \(\nu\) (and not just its restriction to the dense subspace \(\mathcal{K}(X; \mathbf{C}) \otimes_{\mathbf{C}} \mathcal{K}(Y; \mathbf{C})\) of \(\mathcal{K}(X \times Y; \mathbf{C})\) ).

The mapping \((\lambda, \mu) \mapsto \lambda \otimes \mu\) of \(\mathcal{M}(\mathrm{X}; \mathbf{C}) \times \mathcal{M}(\mathrm{Y}; \mathbf{C})\) into \(\mathcal{M}(\mathrm{X} \times \mathrm{Y}; \mathbf{C})\) is obviously bilinear, by virtue of formula (3) of No.~1.

PROPOSITION 1. --- Let \(\lambda\) be a measure on X, \(\mu\) a measure on Y; if \(g \in \mathcal{C}(\mathrm{X};\mathbf{C})\) , \(h \in \mathcal{C}(\mathrm{Y};\mathbf{C})\) , then

\[
(g \cdot \lambda) \otimes (h \cdot \mu) = (g \otimes h) \cdot (\lambda \otimes \mu).\tag{6}
\]

For every function \(f \in \mathcal{K}(\mathrm{X} \times \mathrm{Y}; \mathbf{C})\) , we have, by virtue of formula (3) of No.~1,

\[
\begin{array}{c} \langle f, (g \otimes h) \cdot (\lambda \otimes \mu) \rangle = \int d \lambda (x) \int f (x, y) g (x) h (y) d \mu (y) \\ = \int g (x) d \lambda (x) \int f (x, y) h (y) d \mu (y), \end{array}
\]

which proves formula (6).

PROPOSITION 2. --- The support of the product \(\lambda \otimes \mu\) is equal to the product of the support of \(\lambda\) and the support of \(\mu\) .

We first observe that the relation \(\lambda \otimes \mu = 0\) implies that one of the measures \(\lambda, \mu\) is zero (A, II, §7, No.~7, Prop. 16, (ii)). On the other hand, if U (resp. V) is an open set in X (resp. Y), then the restriction of \(\lambda\otimes\mu\) to the product \(U\times V\) is the product of the restrictions of \(\lambda\) to U and of \(\mu\) to V, as follows from Th. 1 of No.~1 and the definition of the restriction of a measure to an open set (§2, No.~1). It therefore follows that, for the restriction of \(\lambda\otimes\mu\) to \(U\times V\) to be zero, it is necessary and sufficient that either the restriction of \(\lambda\) to U or the restriction of \(\mu\) to V be zero, which proves the proposition, on taking into account the definition of the support of a measure (§2, No.~2).

PROPOSITION 3. --- Let \(\lambda \in \mathcal{M}(\mathrm{X};\mathbf{C})\) , \(\mu \in \mathcal{M}(\mathrm{Y};\mathbf{C})\) . Then

\[
| \lambda \otimes \mu | = | \lambda | \otimes | \mu |.\tag{7}
\]

Let \(f \in \mathcal{K}_{+}(\mathrm{X} \times \mathrm{Y})\) , \(g \in \mathcal{K}(\mathrm{X} \times \mathrm{Y}; \mathbf{C})\) be such that \(|g| \leqslant f\) ; we have (§1, No.~6, formula (13))

\[
\begin{array}{r l} | \langle g, \lambda \otimes \mu \rangle | = & \left| \int d \lambda (x) \int g (x, y)   d \mu (y) \right| \\ & \leqslant \int d | \lambda | (x) \int | g (x, y) |   d | \mu | (y) \\ & = \langle | g |, | \lambda | \otimes | \mu | \rangle \leqslant \langle f, | \lambda | \otimes | \mu | \rangle . \end{array}
\]

It follows that \(\langle f, |\lambda \otimes \mu| \rangle \leqslant \langle f, |\lambda| \otimes |\mu| \rangle\) , and so

\[
| \lambda \otimes \mu | \leqslant | \lambda | \otimes | \mu |.\tag{8}
\]

On the other hand, let \(u \in \mathcal{K}_{+}(\mathrm{X})\) , \(v \in \mathcal{K}_{+}(\mathrm{Y})\) . For every \(\varepsilon > 0\) , there exist \(u_{1} \in \mathcal{K}(\mathrm{X};\mathbf{C})\) , \(v_{1} \in \mathcal{K}(\mathrm{Y};\mathbf{C})\) such that \(|u_{1}| \leqslant u\) , \(|v_{1}| \leqslant v\) and

\[
| \langle u _ {1}, \lambda \rangle | \geqslant \langle u, | \lambda | \rangle - \varepsilon , \quad | \langle v _ {1}, \mu \rangle | \geqslant \langle v, | \mu | \rangle - \varepsilon
\]

(§1, No.~6). It follows that \(|u_1 \otimes v_1| \leqslant u \otimes v\) and

\[
\begin{array}{c} \langle u \otimes v, | \lambda \otimes \mu | \rangle \geqslant | \langle u _ {1} \otimes v _ {1}, \lambda \otimes \mu \rangle | = | \langle u _ {1}, \lambda \rangle \langle v _ {1}, \mu \rangle | \\ \geqslant (\langle u, | \lambda | \rangle - \varepsilon) (\langle v, | \mu | \rangle - \varepsilon). \end{array}
\]

Since \(\varepsilon\) is arbitrary, we infer that

\[
\langle u \otimes v, | \lambda \otimes \mu | \rangle \geqslant \langle u, | \lambda | \rangle \langle v, | \mu | \rangle = \langle u \otimes v, | \lambda | \otimes | \mu | \rangle .
\]

Taking (8) into account, we see that

\[
\langle u \otimes v, | \lambda \otimes \mu | \rangle = \langle u \otimes v, | \lambda | \otimes | \mu | \rangle .
\]

Since every function in \(\mathcal{K}(\mathrm{X};\mathbf{C})\) (resp. \(\mathcal{K}(\mathrm{Y};\mathbf{C})\) ) is a linear combination of functions in \(\mathcal{K}_{+}(\mathrm{X})\) (resp. \(\mathcal{K}_{+}(\mathrm{Y})\) ), the preceding formula remains true for \(u\in\mathcal{K}(\mathrm{X};\mathbf{C})\) and \(v\in\mathcal{K}(\mathrm{Y};\mathbf{C})\) ; the proposition therefore follows from the fact that \(\mathcal{K}(\mathrm{X};\mathbf{C})\otimes_{\mathbf{C}}\mathcal{K}(\mathrm{Y};\mathbf{C})\) is dense in \(\mathcal{K}(\mathrm{X}\times\mathrm{Y};\mathbf{C})\) .

COROLLARY. --- Let \(\lambda \in \mathcal{M}(\mathrm{X};\mathbf{R})\) , \(\mu \in \mathcal{M}(\mathrm{Y};\mathbf{R})\) . Then

\[
\left\{ \begin{array}{l} (\lambda \otimes \mu) ^ {+} = \lambda^ {+} \otimes \mu^ {+} + \lambda^ {-} \otimes \mu^ {-}, \\ (\lambda \otimes \mu) ^ {-} = \lambda^ {+} \otimes \mu^ {-} + \lambda^ {-} \otimes \mu^ {+}. \end{array} \right.\tag{9}
\]

For, by virtue of Prop. 3,

\[
\begin{array}{r l} & {(\lambda \otimes \mu) ^ {+} = \frac {1}{2} \big (\lambda \otimes \mu + | \lambda | \otimes | \mu | \big)} \\ & {\qquad = \frac {1}{2} \big ((\lambda^ {+} - \lambda^ {-}) \otimes (\mu^ {+} - \mu^ {-}) + (\lambda^ {+} + \lambda^ {-}) \otimes (\mu^ {+} + \mu^ {-}) \big)} \\ & {\qquad = \lambda^ {+} \otimes \mu^ {+} + \lambda^ {-} \otimes \mu^ {-}.} \end{array}
\]

The argument for \((\lambda \otimes \mu)^{-}\) is similar.

PROPOSITION 4. --- Let \(\lambda \in \mathcal{M}(\mathrm{X};\mathbf{C})\) , \(\mu \in \mathcal{M}(\mathrm{Y};\mathbf{C})\) . Then

\[
\| \lambda \otimes \mu \| = \| \lambda \| \cdot \| \mu \|,\tag{10}
\]

with the convention that the second member is to be replaced by 0 whenever one of the factors is 0 and the other is +∞. In particular, if λ and μ are bounded then λ⊗μ is bounded.

By the above Proposition 3, and §1, No.~8, Cor. 1 of Prop. 10, we may limit ourselves to the case that \(\lambda\) and \(\mu\) are positive measures. If \(\lambda = 0\) or \(\mu = 0\) , the result is trivial; let us therefore assume that \(\lambda \neq 0\) and \(\mu \neq 0\) . Let us first prove that \(\| \lambda \otimes \mu \| \leqslant \| \lambda \| \cdot \| \mu \|\) . We may suppose \(\lambda\) and \(\mu\) to be bounded. For every \(f \in \mathcal{K}_+(\mathrm{X} \times \mathrm{Y})\) ,

\[
\langle f, \lambda \otimes \mu \rangle = \int d \lambda (x) \int f (x, y) d \mu (y)
\]

and

\[
\int f (x, y) d \mu (y) \leqslant \| f \| \cdot \| \mu \|
\]

for every \(x\in \mathbf{X}\) , therefore

\[
\langle f, \lambda \otimes \mu \rangle \leqslant \| f \| \cdot \| \lambda \| \cdot \| \mu \|,
\]

which proves our assertion. On the other hand, let

\[
\alpha <   \| \lambda \|, \beta <   \| \mu \|
\]

be two real numbers \(\geqslant 0\) . There exist \(g\in \mathcal{K}_{+}(\mathrm{X})\) , \(h\in \mathcal{K}_{+}(\mathrm{Y})\) such that

\[
\| g \| \leqslant 1, \| h \| \leqslant 1, \lambda (g) \geqslant \alpha , \mu (h) \geqslant \beta .
\]

Then \(g \otimes h \in \mathcal{K}_{+}(X \times Y)\) , \(\| g \otimes h \| \leqslant 1\) and \(\langle g \otimes h, \lambda \otimes \mu \rangle \geqslant \alpha \beta\) ; therefore \(\| \lambda \otimes \mu \| \geqslant \alpha \beta\) and finally \(\| \lambda \otimes \mu \| \geqslant \| \lambda \| \cdot \| \mu \|\) , which completes the proof.

